\documentclass[11pt]{article}
\usepackage{amsmath,amssymb,amsthm}
\usepackage[margin=1in]{geometry}
\usepackage{mathtools}
\usepackage[hidelinks]{hyperref}

\newtheorem{theorem}{Theorem}
\newtheorem{lemma}[theorem]{Lemma}
\newtheorem{definition}[theorem]{Definition}
\newtheorem{corollary}[theorem]{Corollary}
\newtheorem{remark}[theorem]{Remark}

\newcommand{\trueT}{\mathbf{t}}
\newcommand{\falseT}{\mathbf{f}}
\newcommand{\code}[1]{\ulcorner #1 \urcorner}
\newcommand{\reify}[1]{\widehat{#1}}
\newcommand{\RecP}{\mathrm{RecP}}
\newcommand{\substOp}{\mathrm{substOp}}
\newcommand{\thmT}{\mathrm{thmT}}
\newcommand{\thFunTFor}{\mathrm{thFunTFor}}
\newcommand{\cstf}{\mathrm{cstf}}
\newcommand{\TreeEq}{\mathrm{TreeEq}}
\newcommand{\IfLf}{\mathrm{IfLf}}
\newcommand{\encHyp}{\mathrm{encHyp}}

\title{V3 Redesign: Closing Two Loopholes in the Guard/Rose
  G\"odel~II Formalization}
\author{}
\date{}

\begin{document}
\maketitle

\begin{abstract}
The V2 formalization of G\"odel~II in the Guard/Rose equational system
(see \texttt{Guard/archive/guard-godelII.tex}) proves a strictly weaker
theorem than the colloquial statement: its object-level proof evaluator
$\thFunTFor$ admits \emph{two} encoding loopholes that make
\emph{every} consistency sentence refutable by pure computation,
independent of the object theory's actual axiomatic content.  This V3
redesign adds a single new primitive, $\RecP : \mathrm{Fun}_2 \to
\mathrm{Fun}_2$, and restructures the evaluator into a
hypothesis-aware $\thmT$, closing both loopholes and producing a
genuine internal G\"odel~II.
\end{abstract}

\section{The two V2 loopholes}

\subsection{\texttt{mkProofEAny} (the encoding loophole)}

V2's Theorem~14 takes \emph{any} equation $H$ (the ambient hypothesis)
and produces a proof-encoding $\code{H}$ such that
$\thFunTFor(\code{H}) = \reify{\mathrm{codeEqn}\,H}$ --- not because
$H$ is actually derivable from itself, but because the encoding
function \texttt{mkProofEAny} constructs a \emph{fake tree} whose
recursive evaluation ignores $H$'s truth and just echoes $H$'s coded
form.  This collapses the constructive statement ``$H \vdash \mathrm{Con}_H
\Rightarrow H \vdash \bot$'' to a pure computation: ANY equation~$H$
satisfies the premise's syntactic shape, so ANY equation refutes
$\mathrm{Con}_H$.

\subsection{\texttt{case19} (the middle-term loophole)}

V2's $\thFunTFor$ handles the \texttt{ruleTrans} case by emitting
$\mathrm{Pair}(\mathrm{left}(sp_1), \mathrm{right}(sp_2))$ without
checking that $\mathrm{right}(sp_1) = \mathrm{left}(sp_2)$.  This
means
$$\mathrm{trans}(\text{refl}\,\trueT, \text{refl}\,\falseT)$$
has $\mathrm{codeEqn}$-evaluation result $\reify{\mathrm{codeEqn}(\trueT = \falseT)}$
--- a fake encoding of the bottom equation built from genuine refl-proofs.
Combined with \texttt{mkProofEAny}, this compounds the first loophole.

\section{The V3 fix}

\subsection{New primitive: \texorpdfstring{$\RecP$}{RecP}}

We extend the functor syntax with
$$
\RecP : \mathrm{Fun}_2 \to \mathrm{Fun}_2,
\qquad
\RecP(s)\,p\,O = O,
\quad
\RecP(s)\,p\,\mathrm{Pair}(a,b)
  = s\,\mathrm{Pair}(p, \mathrm{Pair}(a,b))\,\mathrm{Pair}(\RecP(s)\,p\,a,\;\RecP(s)\,p\,b).
$$
This is a \emph{parameterised} tree-recursion: the step function $s$
receives the RecP-argument $p$ packed alongside the current node
$\mathrm{Pair}(a,b)$, giving closed-form access to $p$ without
embedding it as a free variable inside $s$'s structure.

From $\RecP$ we build
$$
\substOp : \mathrm{Fun}_2
\qquad
\substOp = \RecP(\mathrm{stepSubstP})
$$
with a step function that inspects the current node's tag and, when
it sees $\mathrm{tagVar}$, substitutes according to the parameter pair
$\mathrm{Pair}(t_C, x_C)$.  The load-bearing correctness lemma is
$$
\substOp\,\mathrm{Pair}(\reify{\mathrm{code}\,t},\;\reify{\mathrm{natCode}\,x})\,\reify{\mathrm{code}\,l}
  = \reify{\mathrm{code}(\mathrm{subst}\,x\,t\,l)}.
$$

\subsection{Case 19: middle-term matching}

$$
\mathrm{case19V3} =
\mathrm{Fan}\bigl(\mathrm{Fan}\,\mathrm{recs_{AR}}\,\mathrm{recs_{BL}}\,\TreeEq\bigr)
\;\bigl(\mathrm{Fan}\,(\mathrm{Fan}\,\mathrm{recs_{AL}}\,\mathrm{recs_{BR}}\,\mathrm{Pair})\;(\mathrm{kF_2}\,O)\;\mathrm{Pair}\bigr)\;\IfLf
$$
This inserts a $\TreeEq$ check between $\mathrm{right}(sp_1)$ and
$\mathrm{left}(sp_2)$: on match, emit
$\mathrm{Pair}(\mathrm{left}(sp_1), \mathrm{right}(sp_2))$; on
mismatch, emit the sentinel $O$.  The attempted fake
$\mathrm{trans}(\text{refl}\,\trueT, \text{refl}\,\falseT)$ no longer
typechecks in the V3 system because the middle terms fail to unify.

\subsection{Case 23: dynamic $\mathrm{ruleInst}$ via $\substOp$}

$$
\mathrm{case23V3} =
\mathrm{Fan}\;(\mathrm{Fan}\,\mathrm{origA}\,\mathrm{recs_{BL}}\,\substOp)
\;(\mathrm{Fan}\,\mathrm{origA}\,\mathrm{recs_{BR}}\,\substOp)\;\mathrm{Pair}
$$
V2's \texttt{case23} applied \texttt{substTFor}~(a closed Fun$_1$ with
free variables $v_{11}, v_{12}$ inside it) to each sub-equation side,
and later \emph{ruleInst-bound} $v_{11}, v_{12}$ at proof-construction
time.  V3's \texttt{case23V3} evaluates $\substOp$ at the Fun$_2$-level
with the substitution parameters already packed into the
\texttt{Pair} argument --- no var-capture.

\subsection{Hypothesis check: \texttt{case26}}

At the dispatch bottom we add a new tag $n_{26}$ with
$$
\mathrm{case26}(h) =
\mathrm{Fan}\,(\mathrm{Fan}\,(\mathrm{Lift}\,\mathrm{Snd})\,(\mathrm{kF_2}\,h)\,\TreeEq)\;
(\mathrm{Fan}\,(\mathrm{kF_2}\,h)\,(\mathrm{kF_2}\,O)\,\mathrm{Pair})\;\IfLf
$$
and the $\mathrm{ruleHyp}$ encoding is redefined as $\encHyp(h) =
\mathrm{nd}(\mathrm{natCode}\,n_{26},\;h)$.  The encoding is accepted
only when the body literally equals the ambient hypothesis code $h$,
so \texttt{mkProofEAny} no longer passes:
any encoding of \texttt{ruleHyp} at a fake-hypothesis $H'$ now
evaluates to the sentinel $O$ under $\thmT(\reify{\mathrm{codeEqn}\,H})$.

\section{The V3 G\"odel I argument}

The V2 diagonal $\cstf\,(\reify{\mathrm{templateCode}}) = \reify{cGS}$
translates to V3 via $\substOp$ --- \texttt{diagFTargetV3} in
\texttt{Guard/GodelIV3.agda}.  The key wrinkle: since $\thmT$ is
hypothesis-specific, the V2 trick of substituting a single variable
$v_1$ with $\reify{\mathrm{templateCode}}$ in the template doesn't
suffice: after substitution, the sentence's embedded $\thmT$-hCode
becomes $\reify{\mathrm{templateCode}}$, whereas $\mathrm{thm}_{14\mathrm{E}}\mathrm{V}3$'s
correctness witness produces $\thmT(\reify{cGS})$.

\emph{Fix:} use \emph{two} free variables in the template~---
$v_1$ for the self-code (substituted in the template) and $v_2$ for
$\thmT$'s hCode (kept free until Deriv-level $\mathrm{ruleInst}$).
The G\"odel~I proof chain then threads four $\mathrm{ruleInst}$'s
($v_2 \mapsto \reify{cGS}$, $v_{11}' \mapsto \reify{crTC}$,
$v_{12}' \mapsto \reify{\mathrm{natCode}\,v_1}$, $0 \mapsto enc$) to
bring the sentence's embedded $\thmT$-hCode into alignment with
$\mathrm{corr}\,pe$'s witness.  Both sides of the sentence's $\TreeEq$
collapse to $\reify{cGS}$; \texttt{treeEqSelf} closes the diagonal.

\begin{theorem}[G\"odel~I on V3 --- \texttt{godelIDerivV3} in \texttt{GodelIV3.agda}]
For the V3 G\"odel sentence $\mathrm{gs}_{V3}$:
\[
  \mathrm{gs}_{V3} \vdash \trueT = \falseT.
\]
\end{theorem}

\section{The V3 G\"odel II statement}

\begin{theorem}[G\"odel~II on V3 --- \texttt{godelIIV3} in \texttt{GodelIIV3.agda}]
If $\mathrm{gs}_{V3}$ is consistent, then
$\mathrm{gs}_{V3} \not\vdash \mathrm{Con}_{V3}$,
where
$\mathrm{Con}_{V3} = \TreeEq(\thmT(\reify{cGS})(x),\;\reify{\mathrm{codeEqn}(\trueT = \falseT)}) = \falseT$.
\end{theorem}

The argument: necessitation of \texttt{godelIDerivV3} gives a
\texttt{Prov3}-witness $enc_\bot$ with
$\thmT(\reify{cGS})(enc_\bot) = \reify{\mathrm{codeEqn}(\trueT=\falseT)}$.
Plugging $enc_\bot$ into the consistency sentence via ruleInst and
rewriting the $\thmT$-application using \texttt{corr}, both arguments
of the internal $\TreeEq$ become $\reify{\mathrm{codeEqn}(\trueT=\falseT)}$;
\texttt{treeEqSelf} then contradicts the sentence's claim that this
should equal $\falseT$.

\section{Negative acceptance test}

\texttt{Guard/GodelIITestV3.agda} documents that the V2-style attack
$\mathrm{thm}_{14\mathrm{E}}\mathrm{V}3\mathrm{Trans}(\text{refl}\,\trueT)(\text{refl}\,\falseT)$
is rejected at type-checking time: Agda reports
$\falseT = \mathrm{Pair}(O,O) \neq O = \trueT$ when unifying the two
middle terms.  V3 has no \texttt{mkProofEAny}-like escape hatch, so
the type system structurally rules out the V2 loopholes.

\section{Files and sizes}

All files typecheck under \texttt{--safe --without-K --exact-split}
with the \texttt{agda-2.9.0} binary.  Clean rebuild of
\texttt{Guard/GodelIIV3.agda} (the top of the V3 chain) is $\sim$7~s.

\begin{itemize}
\item \texttt{Guard/Term.agda}, \texttt{Guard/Step.agda} ---
  $\RecP$ primitive and its axioms \texttt{axRecPLf}, \texttt{axRecPNd}.
\item \texttt{Guard/SubstOp.agda} --- $\substOp$ + $\texttt{substOpEquiv}$
  + $\texttt{substOpCorrect}$.
\item \texttt{Guard/ThFunTForV3.agda} --- the hypothesis-aware dispatcher,
  including \texttt{case19V3}, \texttt{case23V3}, \texttt{case26},
  \texttt{case27}, \texttt{case28}.
\item \texttt{Guard/Thm14EV3.agda} --- \texttt{thm14EV3} (total, 29
  Deriv constructors) and the per-rule encodings.
\item \texttt{Guard/SubstTForPrecompV3.agda} --- V3 diagonal machinery
  (templateCode$_{V3}$, cGS$_{V3}$, etc.) plus the fused
  \texttt{fullyInstGS} substEq lemma.
\item \texttt{Guard/GodelIV3.agda} --- \texttt{godelIDerivV3},
  \texttt{diagFTargetV3}.
\item \texttt{Guard/ProvV3.agda} --- \texttt{Prov3},
  \texttt{necessitation}, \texttt{diagContradict}.
\item \texttt{Guard/GodelIIV3.agda} --- \texttt{provBot},
  \texttt{conToBotV3}, \texttt{godelIIV3}.
\item \texttt{Guard/GodelIITestV3.agda} --- negative acceptance test.
\end{itemize}

\end{document}
